\documentclass{article}
\input{../../../lib.sty}

\title{\texttt{fn gamma\_lo\_from\_x}}
\author{Michael Shoemate}

\begin{document}
\maketitle

Proves soundness of \texttt{fn gamma\_lo\_from\_x} in \texttt{mod.rs}.

\section{Hoare Triple}
\subsection*{Precondition}
\texttt{x} is finite.

\subsection*{Pseudocode}
\lstinputlisting[language=Python,firstline=2,escapechar=|]{./pseudocode/gamma_lo_from_x.py}

\subsection*{Postcondition}
\begin{theorem}
\label{postcondition}
For any setting of the input parameter \texttt{x} such that the given precondition holds,
\texttt{gamma\_lo\_from\_x} either returns \texttt{Err(e)},
or returns \texttt{Ok(out)} where
\[
\texttt{out} \le 1 - e^{-\texttt{x}}.
\]
\end{theorem}

\section{Proof}

Assume the precondition is met.

\begin{lemma}
\label{err-e}
\texttt{gamma\_lo\_from\_x} only returns \texttt{Err(e)} if the call to
\texttt{inf\_exp\_m1} on line \ref{line:expm1_hi} returns \texttt{Err(e)}.
\end{lemma}

\begin{proof}
The only fallible operation in \texttt{gamma\_lo\_from\_x} is the invocation of
\texttt{(-x).inf\_exp\_m1()} on line \ref{line:expm1_hi}.
Negation of a finite floating-point value is infallible, and negating the result on the return line
is infallible as well.
Therefore every error returned by \texttt{gamma\_lo\_from\_x} must come from
\texttt{inf\_exp\_m1}.
\end{proof}

\begin{lemma}
\label{ok-out}
If the outcome of \texttt{gamma\_lo\_from\_x} is \texttt{Ok(out)}, then
\[
\texttt{out} \le 1 - e^{-\texttt{x}}.
\]
\end{lemma}

\begin{proof}
Suppose \texttt{gamma\_lo\_from\_x} returns \texttt{Ok(out)}.
Write \texttt{expm1\_hi} for the value returned by \texttt{(-x).inf\_exp\_m1()} on line \ref{line:expm1_hi}.
By the proof definition of \texttt{inf\_exp\_m1},
\[
\texttt{expm1\_hi} \ge e^{-\texttt{x}} - 1.
\]
Multiplying both sides by $-1$ reverses the inequality, giving
\[
-\texttt{expm1\_hi} \le 1 - e^{-\texttt{x}}.
\]
By the return statement of the function, \texttt{out} $= -\texttt{expm1\_hi}$.
Therefore
\[
\texttt{out} \le 1 - e^{-\texttt{x}},
\]
as claimed.
\end{proof}

\begin{proof}[Proof of Theorem~\ref{postcondition}]
Holds by Lemma~\ref{err-e} and Lemma~\ref{ok-out}.
\end{proof}

\end{document}
